%%%%%%%%%%  scan idx 626  |  folio 613  %%%%%%%%%%
En
effet ces conditions sont suffisantes car déjà 4) et 5) entraînent l'effectivité. Réciproquement soit \(M\) inversible sur \(Y\) tel qu'il existe un isomorphisme \(f^*M\simeq L\) compatible avec la donnée de recollement. On peut recouvrir \(Y\) par des ouverts \(V\), au-dessus desquels \(M\) est isomorphe à \({O}_Y\)~; donc la donnée de descente sur \(L\) restreinte au-dessus de \(V=f^{-1}(V)\) est isomorphe à la donnée de descente triviale sur \({O}_X|V\). Enfin les conditions 1)-5) sont de nature locale sur \(Y\) et elles sont vérifiées d'après les hypothèses de 4.1 pour \(L={O}_X\)~; elles sont donc vérifiées pour \(L\).

Il reste donc à représenter successivement chacune des conditions 1) à 5) par des monomorphismes de présentation finie. Pour 1) et 2) on applique le

\textbf{Lemme 4.4.} \uline{Soient} \(S\) \uline{un schéma,} \(\lambda:X\to S\) \uline{et} \(\mu:Y\to S\) \uline{des morphismes propres de présentation finie,} \(f:X\to Y\) \uline{un} \(S\)\uline{-morphisme. Soit} \(F\) \uline{un faisceau sur} \(X\) \uline{de présentation finie et} plat \uline{sur} \(S\) \uline{et, pour un entier} \(n\)\uline{, soit} \(S_n\) \uline{le sous-foncteur de} \(S\) \uline{qui rend} \(R^if_*(F)\) \uline{plat} \uline{sur} \(S\) \uline{pour tout} \(i\ge n\) \uline{(précisément,} \(S_n(T)=\{\varnothing\}\) \uline{si} \(R^i(f_T)_*(F_T)\) \uline{est} \uline{plat} \uline{pour tout} \(i\ge n\), \(S_n(T)=\varnothing\) \uline{sinon). Alors} \(S_n\hookrightarrow S\) \uline{est représentable par un monomorphisme surjectif de présentation finie, et alors pour} \(i\ge n\) \uline{les} \(R^i(f_{S_n})_*(F_{S_n})\) \uline{commutent aux changements de base} \(T\to S_n\)\uline{. De plus} \(S_n\) \uline{coïncide avec le foncteur} \(S'_n\) \uline{qui rend exacts les foncteurs}
\[
  M\longmapsto R^i(f_V)_*(F_V\otimes\lambda_V^*M)\qquad \text{\uline{pour} } i\ge n
\]
\uline{où} \(M\) \uline{parcourt les} \({O}_V\)\uline{-modules quasi-cohérents pour tout ouvert affine} \(V\) \uline{de} \(S\).

En effet l'assertion est locale sur \(S\)~; on peut donc supposer \(S\) quasi-compact. Alors la dimension relative de \(X/Y\) est bornée, soit par \(r\), et pour \(n>r\), \(R^nf_*(F_V\otimes\lambda_V^*M)=0\). L'assertion est donc vraie pour \(n>r\)~;
